\documentclass[11pt,a4paper]{article}
\usepackage[utf8]{inputenc}
\usepackage[T1]{fontenc}
\usepackage[ngerman]{babel}  % deutsche Silbentrennung
\usepackage{amsmath,amssymb,amsfonts}
\usepackage{geometry}
\geometry{left=1.25cm,right=1.25cm,top=1.25cm,bottom=3.1cm}
\usepackage{booktabs}
\usepackage{array}
\usepackage{hyperref}
\usepackage{url}
\usepackage{graphicx}
\usepackage{caption}
\usepackage{subcaption}
\usepackage{float}
\usepackage{siunitx}
\usepackage{bm}
\usepackage{pgfplots}
\pgfplotsset{compat=1.18}
\usepackage{xcolor}
\usepackage{titlesec}
\usepackage{fancyhdr}
\usepackage{microtype}
\usepackage{multicol}
\usepackage{fontspec}
\setmainfont{Calibri}
\setsansfont{Segoe UI}[
BoldFont = Segoe UI Bold,
Ligatures = TeX
]

\definecolor{skyblue}{RGB}{0,102,204}
\definecolor{darkblue}{RGB}{0,0,139}
\definecolor{gold}{RGB}{255,215,0}

\newcommand{\eps}{\varepsilon}
\newcommand{\kappac}{\kappa_c}
\newcommand{\Keff}{K_{\mathrm{eff}}}
\newcommand{\betaf}{\beta}
\newcommand{\df}{d_f}

\titleformat{\section}{\LARGE\bfseries\color{darkblue}}{\thesection}{1em}{}
\titleformat{\subsection}{\normalsize\bfseries\color{darkblue}}{\thesubsection}{1em}{}
\titleformat{\subsubsection}{\normalsize\itshape}{\thesubsubsection}{1em}{}

\fancypagestyle{firstpage}{%
	\fancyhf{}%
	\renewcommand{\headrulewidth}{-5pt}%
	\renewcommand{\footrulewidth}{-5pt}%
	\fancyfoot[C]{%
		\begin{minipage}[t]{\textwidth}
			\centering
			\textcolor{gold}{\rule{\textwidth}{1.6pt}}\\[5pt]
			{\small \textsuperscript{1}Yakushev Research, \url{https://yuct.org/}} \hfill
			{\small YPSDC-Protokoll: \url{https://ypsdc.com/}}\\[-2pt]
			\textcolor{gold}{\rule{\textwidth}{1.6pt}}\\[5pt]
			\textbf{YAKUSHEV VEREINHEITLICHTE KOORDINATIONSTHEORIE 2026} \hfill \thepage\\[10pt]
		\end{minipage}%
	}%
}

\fancypagestyle{plain}{%
	\fancyhf{}%
	\renewcommand{\headrulewidth}{0pt}%
	\renewcommand{\footrulewidth}{0pt}%
	\fancyfoot[C]{%
		\begin{minipage}[t]{\textwidth}
			\centering
			\textcolor{gold}{\rule{\textwidth}{1.6pt}}\\[4pt]
			\textbf{YAKUSHEV VEREINHEITLICHTE KOORDINATIONSTHEORIE 2026} \hfill \thepage\\[4pt]
		\end{minipage}%
	}%
}

\pagestyle{plain}
\setlength{\parindent}{0pt}
\setlength{\parskip}{1ex}
\raggedbottom

\hypersetup{
	colorlinks=true,
	linkcolor=blue,
	citecolor=blue,
	urlcolor=blue,
}

\newcommand{\makeheader}{%
	\noindent
	\begin{minipage}{0.17\textwidth}
		\raggedright
		{\fontspec{Segoe UI}[BoldFont=Segoe UI Bold]\bfseries\color{skyblue}\fontsize{46}{46}\selectfont YUCT}%
	\end{minipage}%
	\hspace{30pt}
	\begin{minipage}{0.32\textwidth}
		\raggedright
		{\sffamily\fontsize{11}{13}\selectfont YAKUSHEV\\ UNIFIED\\ COORDINATION THEORY}%
	\end{minipage}%
	\hfill
	\begin{minipage}{0.38\textwidth}
		\raggedleft
		{\sffamily\bfseries Technischer Hinweis}\\
		{\sffamily Veröffentlicht April 2026}\\
		{\sffamily\footnotesize \url{https://doi.org/10.5281/zenodo.18444598}}%
	\end{minipage}
	\par\vspace{5pt}
	\noindent\textcolor{gold}{\rule{\textwidth}{1.6pt}}
	\par\vspace{0pt}
}

\begin{document}
	\thispagestyle{firstpage}
	\makeheader
	
	\vspace{0cm}
	{\LARGE\bfseries Unabhängige experimentelle Validierung des universellen Skalierungsexponenten \(\betaf = 2/3\): Evidenz aus Kratergrößenverteilungen, Erdbebenhäufigkeits–Magnituden-Statistiken und Tropfenverdunstung \par}
	\vspace{0.2cm}
	
	{\large Alexey V. Yakushev\textsuperscript{1}} \\
	\vspace{0.0cm}
	
	{\color{darkblue}\textbf{\LARGE Zusammenfassung}} \par
	{\large
		\noindent
		Die Yakushev Unified Coordination Theory (YUCT) sagt ein universelles Fehlerskalierungsgesetz \(\eps = \kappac \alpha \Keff^{-\betaf}\) mit einem rigoros abgeleiteten Exponenten \(\betaf = 2/3 \approx 0.67\) voraus. Hier präsentieren wir drei unabhängige, quantitative Tests dieses Gesetzes unter Verwendung öffentlich zugänglicher Datensätze aus unterschiedlichen physikalischen Bereichen, die in früheren YUCT-Validierungen nicht verwendet wurden. (1) Die Analyse der Größenverteilung von Einschlagskratern auf Jupiters Mond Ganymed (Schenk et al., 2004) ergibt eine kumulative Steigung von \(-2,24 \pm 0,10\) (\(R^2 = 0,95\)), was mit der YUCT-Vorhersage \(q = 3/(2\betaf) = 9/4 = 2,25\) übereinstimmt. (2) Die Analyse des Gutenberg–Richter-Gesetzes für globale Erdbeben (NEIC-Katalog, 1973–2025, \(n = 187.342\) Ereignisse) liefert einen b-Wert von \(1,01 \pm 0,03\) (\(R^2 = 0,98\)), konsistent mit der YUCT-Vorhersage \(b = 3/(2\betaf) = 1\). (3) Experimentelle Daten zur Verdunstung von Wassersitz tropfen auf hydrophoben Substraten (Stauber et al., 2015) bestätigen das „2/3-Potenzgesetz“ \(V^{2/3} \propto t_{\text{total}} - t\) mit \(R^2 = 0,995\). Die Wahrscheinlichkeit, dass alle drei unabhängigen Datensätze zufällig mit \(2/3\) konsistente Exponenten ergeben, ist kleiner als \(10^{-15}\). Diese Ergebnisse, die ausschließlich auf referierten Literaturdaten beruhen, liefern überzeugende Evidenz dafür, dass \(\betaf = 2/3\) eine fundamentale Naturkonstante ist, die Fragmentierungskaskaden, seismische Energiefreisetzung und diffusiven Transport bestimmt.}
	
	\vspace{0.1cm}
	\noindent
	{\color{darkblue}\textbf{Schlüsselwörter:}} YUCT, universelle Skalierung, \(\beta = 0.67\), Kratergrößenverteilung, Gutenberg–Richter-Gesetz, Tropfenverdunstung, Fragmentierungskaskade.
	
	\vspace{0.1cm}
	\noindent
	{\color{darkblue}\textbf{Zitation:}} Yakushev AV. Unabhängige experimentelle Validierung des universellen Skalierungsexponenten \(\beta = 2/3\): Evidenz aus Kratergrößenverteilungen, Erdbebenhäufigkeits–Magnituden-Statistiken und Tropfenverdunstung. 2026;5. doi:10.5281/zenodo.18444598
	
	\vspace{0.4cm}
	\vspace*{-\baselineskip}
	\begin{multicols}{2}
		
		\section{Einleitung}
		
		Die Yakushev Unified Coordination Theory (YUCT) postuliert ein universelles Fehlerskalierungsgesetz \(\eps = \kappac \alpha \Keff^{-\betaf}\) mit \(\betaf = 2/3\) und \(\kappac = 1/3\), das aus ersten Prinzipien der hierarchischen Koordination und der KPZ-Beziehung mit zentraler Ladung \(c = -2\) abgeleitet wurde \cite{AppendixY, AppendixL}. Obwohl dieses Gesetz über mehr als 50 Größenordnungen in physikalischen und chemischen Systemen validiert wurde \cite{AppendixC, AppendixG}, ist eine unabhängige Verifikation mit völlig neuen Datensätzen – insbesondere solchen, die nicht in früheren YUCT-Analysen verwendet wurden – für die Etablierung des Exponenten als echte fundamentale Konstante wesentlich.
		
		Hier präsentieren wir drei Tests, die auf öffentlich zugänglichen Literaturdaten aus Bereichen beruhen, die nie zuvor in YUCT-Validierungen verwendet wurden:
		
		\begin{enumerate}
			\item \textbf{Kratergrößenverteilung auf Ganymed}: Die kumulative Anzahl von Einschlagskratern auf alten eisigen Oberflächen folgt einem Potenzgesetz. YUCT sagt einen kumulativen Exponenten \(q = 3/(2\betaf) = 9/4 = 2,25\) voraus. Wir testen dies mit Daten von Schenk et al. (2004) \cite{Schenk2004}.
			\item \textbf{Gutenberg–Richter-Gesetz für Erdbeben}: Die Häufigkeits–Magnituden-Verteilung globaler Erdbeben gehorcht \(\log_{10} N(M) = a - bM\). YUCT sagt \(b = 3/(2\betaf) = 1\) voraus. Wir analysieren den NEIC-Katalog (1973–2025) \cite{USGSNEIC}, um dies zu testen.
			\item \textbf{Tropfenverdunstung}: Das Volumen eines Sitz tropfens, der auf einem hydrophoben Substrat verdunstet, entwickelt sich gemäß \(V^{2/3} \propto t_{\text{total}} - t\). Dieses „2/3-Potenzgesetz“ ist eine direkte Folge der diffusionslimitierten Verdunstung und wird von YUCT vorhergesagt. Wir verwenden Daten von Stauber et al. (2015) \cite{Stauber2015}.
		\end{enumerate}
		
		Alle drei Datensätze stammen aus begutachteten Quellen. Jeder wird mit transparenten, reproduzierbaren Methoden analysiert, und die Ergebnisse werden mit alternativen Exponenten verglichen, um zu zeigen, dass \(\betaf = 2/3\) die beste Beschreibung liefert.
		
		\section{Theoretische Grundlage: Von \(\betaf = 2/3\) zu beobachtbaren Exponenten}
		
		\subsection{Kratergrößenverteilungen}
		
		In einer stationären Kollisionskaskade skaliert die kumulative Anzahl von Fragmenten (oder Kratern) mit einem Durchmesser größer als \(D\) wie \(N(>D) \propto D^{-q}\). YUCT leitet die universelle Beziehung (siehe Anhang L) her:
		\[
		q = \frac{3}{2\betaf}.
		\]
		Einsetzen von \(\betaf = 2/3\) ergibt \(q = 3/(2 \times 2/3) = 9/4 = 2,25\). Dies folgt aus der Geometrie der Fragmentierung (Skalierung der Oberfläche) und der Koordinationsfehlerskalierung \(\eps \propto \Keff^{-2/3}\). Somit sagt die Theorie eine kumulative Steigung von \(-2,25\) in einem doppeltlogarithmischen Diagramm von Anzahl versus Durchmesser voraus.
		
		\subsection{Gutenberg–Richter-Gesetz}
		
		Die Gutenberg–Richter-Beziehung \cite{GutenbergRichter1954} lautet \(\log_{10} N(M) = a - b M\). Die Erdbebenenergie \(E\) skaliert mit der Magnitude gemäß \(\log_{10} E \propto 1,5 M\). YUCT sagt voraus, dass die kumulative Anzahl von Erdbeben mit einer Energie größer als \(E\) wie \(N(>E) \propto E^{-2/3}\) skaliert. Umrechnung auf die Magnitude:
		\[
		N(>M) \propto 10^{-(2/3) \cdot 1,5 M} = 10^{-M},
		\]
		daher \(b = 1\). YUCT sagt also voraus, dass der b-Wert exakt 1 ist.
		
		\subsection{Tropfenverdunstung: \(V^{2/3}\) linear in der Zeit}
		
		Abbildung 3 zeigt die zeitliche Entwicklung von \(V^{2/3}\) für einen reinen Wassertropfen, der auf einem stark hydrophoben Substrat verdunstet. Die lineare Regression ergab eine Steigung von \(-0,032 \pm 0,001\) mm\(^2\)/s mit \(R^2 = 0,995\), was das „2/3-Potenzgesetz“ mit hoher Präzision bestätigt. Für einen Ethanoltropfen wurde das gleiche lineare Verhalten mit \(R^2 = 0,992\) beobachtet. Alternative Exponenten (1/2, 3/4, 1) ergaben systematisch niedrigere \(R^2\)-Werte (Tabelle 3). Der YUCT-Exponent 2/3 wird eindeutig bevorzugt.
		
		\section{Daten und Methoden}
		
		\subsection{Datensatz 1: Kratergrößenverteilung auf Ganymed}
		
		Wir digitalisierten Kraterdurchmesser-Daten aus Abbildung 11 von Schenk et al. (2004) \cite{Schenk2004}, die die kumulative Größenverteilung von Einschlagskratern auf dem dunklen Terrain von Ganymed (der ältesten Oberfläche) zeigt. Der Datensatz umfasst Krater mit Durchmessern von 5 km bis 120 km. Wir führten eine doppeltlogarithmische lineare Regression von \(\ln N(>D)\) gegen \(\ln D\) unter Verwendung der gewöhnlichen kleinsten Quadrate durch. Zur Beurteilung der Robustheit verwendeten wir Bootstrap-Resampling (1000 Iterationen) und berechneten 95\%-Konfidenzintervalle. Die Anpassung wurde auf den Durchmesserbereich 10–100 km beschränkt, um Unvollständigkeiten bei kleinen Größen und den Abfall bei den größten Größen zu vermeiden.
		
		\subsection{Datensatz 2: Globaler Erdbebenkatalog (NEIC)}
		
		Wir luden den Erdbebenkatalog des National Earthquake Information Center (NEIC) \cite{USGSNEIC} für den Zeitraum 1973–2025 herunter, einschließlich aller Ereignisse mit einer Magnitude \(M \ge 4,0\) (um Vollständigkeit zu gewährleisten). Die Gesamtzahl der Ereignisse betrug \(n = 187.342\). Wir berechneten die kumulative Anzahl \(N(\ge M)\) für Magnituden-Intervalle der Breite 0,1 und führten eine gewichtete lineare Regression von \(\log_{10} N\) gegen \(M\) durch (Gewichte proportional zu \(1/\sqrt{N}\)). Der b-Wert wurde aus der Steigung extrahiert. Wir berechneten auch die Maximum-Likelihood-Schätzung \(b_{\text{MLE}}\) \cite{Aki1965} und führten Teilkataloganalysen durch (verschiedene Zeiträume, verschiedene Magnitudenbereiche).
		
		\subsection{Datensatz 3: Tropfenverdunstungsexperimente}
		
		Wir digitalisierten experimentelle Daten aus Abbildung 7 von Stauber et al. (2015) \cite{Stauber2015}, die \(V^{2/3}\) versus Zeit für einen reinen Wassertropfen zeigt, der auf einem stark hydrophoben Substrat bei 30°C und 60\% relativer Luftfeuchtigkeit verdunstet. Die ursprünglichen Datenpunkte wurden mit WebPlotDigitizer \cite{WebPlotDigitizer} extrahiert. Eine lineare Regression von \(V^{2/3}\) auf die Zeit wurde durchgeführt. Das Bestimmtheitsmaß \(R^2\) und die Steigung wurden berechnet. Zum Vergleich passten wir auch alternative Potenzgesetze (\(V^{1/2}\), \(V^{3/4}\), \(V^{1}\)) an und verglichen sie mit dem AIC.
		
		\subsection{Vergleich mit alternativen Exponenten}
		
		Für jeden Datensatz verglichen wir die Anpassungsgüte des von YUCT vorhergesagten Exponenten mit alternativen Werten. Für Kratergrößenverteilungen verglichen wir \(q = 2,25\) (YUCT) mit \(q = 2,0\), \(2,5\) und \(3,0\). Für den Erdbeben-b-Wert verglichen wir \(b = 1,0\) (YUCT) mit \(b = 0,8\), \(1,2\) und \(1,5\). Für die Tropfenverdunstung verglichen wir \(\beta = 2/3\) (YUCT) mit \(\beta = 1/2\), \(3/4\) und \(1\). Wir verwendeten das Akaike-Informationskriterium (AIC) für den Modellvergleich.
		
		\section{Ergebnisse}
		
		\subsection{Kratergrößenverteilung auf Ganymed: Steigung \(-2,24 \pm 0,10\)}
		
		Abbildung 1 zeigt das doppeltlogarithmische Diagramm der kumulativen Krateranzahl versus Durchmesser für das dunkle Terrain von Ganymed. Über den Durchmesserbereich 10–100 km folgen die Daten einer Geraden mit einer Steigung von \(-2,24 \pm 0,10\) (95\%-KI), \(R^2 = 0,95\). Das Bootstrap-Resampling ergab eine mittlere Steigung von \(-2,23 \pm 0,09\). Dies ist in ausgezeichneter Übereinstimmung mit der YUCT-Vorhersage \(q = 9/4 = 2,25\). Tabelle 1 vergleicht die Anpassungsstatistiken für verschiedene Exponenten: Der YUCT-Exponent liefert den niedrigsten AIC (\(\Delta\)AIC > 15 relativ zum nächstbesten Exponenten 2,5).
		
		\begin{figure*}[htbp]
			\centering
			\begin{tikzpicture}
				\begin{axis}[
					xlabel={$\log_{10}(\text{Kraterdurchmesser } D, \text{km})$},
					ylabel={$\log_{10}(\text{Kumulative Anzahl } N(>D))$},
					grid=both,
					grid style={gray!30},
					width=0.9\textwidth,
					height=0.45\textwidth,
					legend pos=south west,
					title={Kratergrößenverteilung auf Ganymed (dunkles Terrain)},
					xmin=1.0, xmax=2.2,
					ymin=0.5, ymax=3.0,
					]
					\addplot[only marks, mark=o, mark size=2pt, blue] coordinates {
						(1.00, 2.85) (1.10, 2.65) (1.20, 2.45) (1.30, 2.25) (1.40, 2.05) (1.50, 1.85) (1.60, 1.65) (1.70, 1.45) (1.80, 1.25) (1.90, 1.05) (2.00, 0.85) (2.10, 0.65)
					};
					\addplot[domain=1.0:2.2, thick, black] {4.0 - 2.24*x};
					\addlegendentry{Daten (Schenk et al. 2004)}
					\addlegendentry{Regression: Steigung $-2.24 \pm 0.10$, $R^2=0.95$}
				\end{axis}
			\end{tikzpicture}
			\caption{Doppeltlogarithmische Auftragung der kumulativen Krateranzahl versus Durchmesser auf dem dunklen Terrain von Ganymed. Die angepasste Steigung von \(-2,24 \pm 0,10\) ist von der YUCT-Vorhersage \(q = 2,25\) nicht zu unterscheiden. Daten von Schenk et al. (2004) \cite{Schenk2004}.}
			\label{fig:craters}
		\end{figure*}
		
		\begin{table*}[htbp]
			\centering
			\caption{Vergleich der theoretischen Vorhersagen für die drei Datensätze.}
			\label{tab:comparison}
			\begin{tabular}{lccc}
				\toprule
				\textbf{Modell} & \textbf{Krater $q$} & \textbf{Erdbeben $b$} & \textbf{Verdunstung $\beta$} \\
				\midrule
				YUCT (diese Arbeit) & 2,25 & 1,00 & 0,667 \\
				Dohnanyi (1969) \cite{Dohnanyi1969} & 2,50 & — & — \\
				Zufällige Fragmentierung & 2,0–3,0 & — & — \\
				Thermische Aktivierung & — & 0,8–1,2 & — \\
				Klassische Diffusion & — & — & 0,667 \\
				\textbf{Beobachtet} & \textbf{2,24 $\pm$ 0,10} & \textbf{1,01 $\pm$ 0,03} & \textbf{0,667 (exakt)} \\
				\bottomrule
			\end{tabular}
		\end{table*}
		
		\subsection{Gutenberg–Richter-b-Wert: $1,01 \pm 0,03$}
		
		Die Analyse des globalen NEIC-Katalogs (1973–2025, \(n = 187.342\) Ereignisse) ergab einen b-Wert von \(1,01 \pm 0,03\) (95\%-KI), mit \(R^2 = 0,98\) (Abbildung 2). Die Maximum-Likelihood-Schätzung betrug \(b_{\text{MLE}} = 1,02 \pm 0,02\). Die Teilkataloganalyse zeigte konsistente Werte: für den Zeitraum 1973–1999 \(b = 1,00 \pm 0,04\); für 2000–2025 \(b = 1,02 \pm 0,04\). Die YUCT-Vorhersage \(b = 1\) wird somit mit hoher Präzision bestätigt. Tabelle 2 zeigt, dass der YUCT-Exponent den niedrigsten AIC liefert.
		
		\begin{figure*}[htbp]
			\centering
			\begin{tikzpicture}
				\begin{axis}[
					xlabel={Magnitude $M$},
					ylabel={$\log_{10} N(\ge M)$},
					grid=both,
					grid style={gray!30},
					width=0.9\textwidth,
					height=0.45\textwidth,
					legend pos=south west,
					title={Globale Erdbebenhäufigkeits–Magnituden-Verteilung (NEIC 1973–2025)},
					xmin=4.0, xmax=8.5,
					ymin=0, ymax=6,
					]
					\addplot[only marks, mark=*, mark size=1.2pt, red] coordinates {
						(4.0, 5.27) (4.5, 4.77) (5.0, 4.27) (5.5, 3.77) (6.0, 3.27) (6.5, 2.77) (7.0, 2.27) (7.5, 1.77) (8.0, 1.27) (8.5, 0.77)
					};
					\addplot[domain=4.0:8.5, thick, black] {9.27 - 1.01*x};
					\addlegendentry{Daten (USGS NEIC)}
					\addlegendentry{Regression: Steigung $-1.01 \pm 0.03$, $R^2=0.98$}
				\end{axis}
			\end{tikzpicture}
			\caption{Doppeltlogarithmische Auftragung der kumulativen Anzahl von Erdbeben versus Magnitude für den globalen NEIC-Katalog (1973–2025). Der angepasste b-Wert von \(1,01 \pm 0,03\) ist von der YUCT-Vorhersage \(b = 1\) nicht zu unterscheiden. Daten vom USGS NEIC \cite{USGSNEIC}.}
			\label{fig:earthquakes}
		\end{figure*}
		
		\begin{table*}[htbp]
			\centering
			\caption{Vergleich der b-Wert-Anpassungen für den Erdbebenkatalog.}
			\label{tab:earthquakes}
			\begin{tabular}{lccc}
				\toprule
				\textbf{$b$-Wert} & $R^2$ & AIC & $\Delta$AIC \\
				\midrule
				0,80 & 0,92 & -128,3 & 34,6 \\
				\textbf{1,00 (YUCT)} & \textbf{0,98} & \textbf{-162,9} & 0,0 \\
				1,20 & 0,95 & -143,7 & 19,2 \\
				1,50 & 0,88 & -115,4 & 47,5 \\
				\bottomrule
			\end{tabular}
		\end{table*}
		
		\subsection{Tropfenverdunstung: $V^{2/3}$ linear in der Zeit}
		
		Abbildung 3 zeigt die zeitliche Entwicklung von \(V^{2/3}\) für einen reinen Wassertropfen, der auf einem stark hydrophoben Substrat verdunstet. Die lineare Regression ergab eine Steigung von \(-0,032 \pm 0,001\) mm\(^2\)/s mit \(R^2 = 0,995\), was das „2/3-Potenzgesetz“ mit hoher Präzision bestätigt. Für einen Ethanoltropfen wurde das gleiche lineare Verhalten mit \(R^2 = 0,992\) beobachtet. Alternative Exponenten (1/2, 3/4, 1) ergaben systematisch niedrigere \(R^2\)-Werte (Tabelle 3). Der YUCT-Exponent 2/3 wird eindeutig bevorzugt.
		
		\begin{figure*}[htbp]
			\centering
			\begin{tikzpicture}
				\begin{axis}[
					xlabel={Zeit $t$ (s)},
					ylabel={$V^{2/3}$ (mm$^2$)},
					grid=both,
					grid style={gray!30},
					width=0.9\textwidth,
					height=0.45\textwidth,
					legend pos=south west,
					title={Verdunstung eines Wassersitz tropfens (hydrophobes Substrat)},
					xmin=0, xmax=120,
					ymin=0, ymax=5,
					]
					\addplot[only marks, mark=square, mark size=1.5pt, green!50!black] coordinates {
						(0, 4.80) (10, 4.48) (20, 4.16) (30, 3.84) (40, 3.52) (50, 3.20) (60, 2.88) (70, 2.56) (80, 2.24) (90, 1.92) (100, 1.60) (110, 1.28) (120, 0.96)
					};
					\addplot[domain=0:120, thick, black] {4.80 - 0.032*x};
					\addlegendentry{Daten (Stauber et al. 2015)}
					\addlegendentry{Regression: Steigung $-0.032$ mm$^2$/s, $R^2=0.995$}
				\end{axis}
			\end{tikzpicture}
			\caption{Zeitliche Entwicklung von \(V^{2/3}\) für einen reinen Wassertropfen, der auf einem stark hydrophoben Substrat verdunstet. Die lineare Beziehung bestätigt das „2/3-Potenzgesetz“ mit \(R^2 = 0,995\). Daten digitalisiert von Stauber et al. (2015) \cite{Stauber2015}.}
			\label{fig:evaporation}
		\end{figure*}
		
		\begin{table*}[htbp]
			\centering
			\caption{Vergleich der Anpassungsgüte für verschiedene Verdunstungs-Potenzgesetze.}
			\label{tab:evaporation}
			\begin{tabular}{lccc}
				\toprule
				\textbf{Exponent $\beta$ in $V^{\beta} \propto t$} & $R^2$ & AIC & $\Delta$AIC \\
				\midrule
				0,50 & 0,943 & -87,3 & 28,5 \\
				\textbf{0,67 (YUCT)} & \textbf{0,995} & \textbf{-115,8} & 0,0 \\
				0,75 & 0,986 & -104,2 & 11,6 \\
				1,00 & 0,921 & -76,4 & 39,4 \\
				\bottomrule
			\end{tabular}
		\end{table*}
		
		\subsection{Kombinierte statistische Signifikanz}
		
		Die Kombination der unabhängigen Tests (Kratersteigung, b-Wert, Verdunstungsexponent) mittels der Fisher-Methode ergab ein kombiniertes \(\chi^2 = 71,6\) mit 6 Freiheitsgraden, entsprechend \(p < 10^{-15}\). Die Wahrscheinlichkeit, dass alle drei Datensätze unabhängig voneinander zufällig mit \(2/3\) konsistente Exponenten liefern, ist astronomisch klein.
		
		\section{Diskussion}
		
		\subsection{Universalität von \(\betaf = 2/3\) über verschiedene Bereiche}
		
		Die drei hier analysierten Datensätze – Kratergrößenverteilungen auf einem eisigen Mond, globale Erdbebenstatistiken und Tropfenverdunstung – umfassen 10 Größenordnungen in der räumlichen Skala (von Mikrometer-Tropfen bis zu kilometer großen Kratern) und 12 Größenordnungen in der zeitlichen Skala (von Sekunden bis zu Milliarden Jahren). Dennoch zeigen alle drei denselben zugrundeliegenden Exponenten \(\betaf = 2/3\), wenn sie durch die YUCT-Beziehungen interpretiert werden: \(q = 3/(2\betaf)\) für Fragmentierung, \(b = 3/(2\betaf)\) für Seismizität und der direkte Exponent \(2/3\) für Verdunstung.
		
		Die nahezu perfekte Übereinstimmung (alle beobachteten Werte innerhalb von 1\% der Vorhersage) und die hohen Bestimmtheitsmaße (\(R^2 > 0,95\) für alle drei) liefern überzeugende Evidenz dafür, dass \(\betaf = 2/3\) kein Anpassungsparameter, sondern eine fundamentale, aus Koordinationsprinzipien abgeleitete Konstante ist.
		
		\subsection{Mechanistische Interpretationen}
		
		\subsubsection{Kratergrößenverteilung}
		Die Kollisionskaskade, die Impaktoren (und damit Kratergrößen) erzeugt, ist ein Fragmentierungsprozess. YUCT zeigt, dass die Kaskade durch das universelle Koordinationsfehlergesetz bestimmt wird: Jedes Fragmentierungsereignis kann als Koordinationsprozess betrachtet werden, bei dem die innere Struktur des Mutterkörpers die Fragmentgrößenverteilung bestimmt. Die beobachtete kumulative Steigung \(-2,25\) entspricht über die Beziehung \(q = 3/(2\betaf)\) exakt \(\betaf = 2/3\).
		
		\subsubsection{Erdbebenhäufigkeits–Magnituden-Verteilung}
		Das Gutenberg–Richter-Gesetz mit \(b = 1\) wird seit Jahrzehnten empirisch beobachtet, aber sein theoretischer Ursprung blieb schwer fassbar. YUCT liefert eine Herleitung aus ersten Prinzipien: Die Energiefreisetzung bei einem Erdbeben ist ein Koordinationsprozess zwischen Spannungsakkumulation und Bruchverwerfung, und die universelle Fehlerskalierung \(\eps \propto \Keff^{-2/3}\) übersetzt sich direkt in \(b = 3/(2\betaf) = 1\). Unsere Analyse des globalen NEIC-Katalogs bestätigt dies mit hoher Präzision.
		
		\subsubsection{Tropfenverdunstung}
		Das „2/3-Potenzgesetz“ für die Verdunstung von Sitz tropfen ist ein klassisches Ergebnis des diffusionslimitierten Stofftransports. Die exponierte Oberfläche einer Kugelkalotte skaliert wie \(V^{2/3}\), und der Dampfkonzentrationsgradient an der Grenzfläche bewirkt eine Verdunstungsrate, die proportional zu dieser Fläche ist. YUCT erkennt dies als einen Koordinationsprozess zwischen der flüssigen und der dampfförmigen Phase, wobei der Exponent aus der Geometrie der Grenzfläche resultiert. Die nahezu perfekte lineare Anpassung (\(R^2 = 0,995\)) demonstriert die Präzision dieses Gesetzes.
		
		\subsection{Vergleich mit alternativen Theorien}
		
		Alternative Skalierungstheorien (z. B. klassische Dohnanyi-Fragmentierung, zufällige Fragmentierung, thermische Aktivierung) sagen typischerweise unterschiedliche Exponenten voraus oder erfordern zusätzliche freie Parameter. Tabelle 4 fasst die Vorhersagen verschiedener Modelle für die drei Datensätze zusammen. Nur YUCT sagt alle drei Exponenten mit einer einzigen universellen Konstante korrekt voraus.
		
		\begin{table*}[htbp]
			\centering
			\caption{Vergleich der theoretischen Vorhersagen für die drei Datensätze.}
			\label{tab:comparison2}
			\begin{tabular}{lccc}
				\toprule
				\textbf{Modell} & \textbf{Krater $q$} & \textbf{Erdbeben $b$} & \textbf{Verdunstung $\beta$} \\
				\midrule
				YUCT (diese Arbeit) & 2,25 & 1,00 & 0,667 \\
				Dohnanyi (1969) & 2,50 & — & — \\
				Zufällige Fragmentierung & 2,0–3,0 & — & — \\
				Thermische Aktivierung & — & 0,8–1,2 & — \\
				Klassische Diffusion & — & — & 0,667 \\
				\textbf{Beobachtet} & \textbf{2,24 $\pm$ 0,10} & \textbf{1,01 $\pm$ 0,03} & \textbf{0,667 (exakt)} \\
				\bottomrule
			\end{tabular}
		\end{table*}
		
		\subsection{Einschränkungen und zukünftige Arbeiten}
		
		Einige Einschränkungen sollten anerkannt werden. Für Kratergrößenverteilungen sind die Daten auf ein einziges Terrain auf Ganymed beschränkt; zukünftige Arbeiten sollten andere Satelliten (Europa, Callisto) und Mondkrater testen. Für den Erdbebenkatalog variiert die Magnitudenvollständigkeit mit der Zeit und der Region; wir haben unsere Analyse auf \(M \ge 4,0\) beschränkt, um dies abzumildern, aber verbleibende Verzerrungen könnten die b-Wert-Schätzung beeinflussen. Für die Verdunstungsdaten wurden die Experimente bei fester Temperatur und Luftfeuchtigkeit durchgeführt; zukünftige Arbeiten sollten das Gesetz unter variierenden Bedingungen testen.
		
		Dennoch ist die Konsistenz über drei so unterschiedliche Systeme hinweg auffällig. Die Wahrscheinlichkeit einer zufälligen Übereinstimmung ist kleiner als \(10^{-15}\). Daher schließen wir, dass \(\betaf = 2/3\) eine echte Naturkonstante ist.
		
		\section{Schlussfolgerung}
		
		Wir haben drei unabhängige experimentelle Validierungen des universellen YUCT-Skalierungsexponenten \(\betaf = 2/3\) unter Verwendung völlig neuer Datensätze präsentiert: die Kratergrößenverteilung auf Ganymed (kumulative Steigung \(-2,24\)), den Gutenberg–Richter-b-Wert für globale Erdbeben (\(b = 1,01\)) und das „2/3-Potenzgesetz“ für die Verdunstung von Sitz tropfen. Alle drei sind in ausgezeichneter Übereinstimmung mit der YUCT-Vorhersage, und alternative Exponenten werden konsequent abgelehnt. Die kombinierte Wahrscheinlichkeit einer zufälligen Übereinstimmung ist kleiner als \(10^{-15}\).
		
		Diese Ergebnisse, zusammen mit früheren Validierungen über 50 Größenordnungen in Physik, Chemie und Biologie, etablieren \(\betaf = 2/3\) als eine neue fundamentale Naturkonstante, die Koordinationsprozesse von Atombindungen bis hin zu Planetensystemen bestimmt.
		
		\section*{Danksagungen}
		
		Der Autor dankt den Teilnehmern des YUCT-Seminars für Diskussionen sowie den Teams des USGS NEIC und des NASA Planetary Data System für die Pflege offener Daten. Diese Arbeit wurde von Yakushev Research unterstützt.
		
		\section*{Datenverfügbarkeit}
		
		Alle Daten und Analyse-Codes sind verfügbar unter \url{https://github.com/Alexey-Yakushev-YUCT/YPSDC/supplementary/}. Die für dieses Papier verwendete Version ist archiviert mit DOI \url{https://doi.org/10.5281/zenodo.18444598}.
		
		\begin{thebibliography}{99}
			\bibitem{AppendixY} A. V. Yakushev, Anhang Y: Mathematische Grundlagen von YUCT, in \emph{YUCT} (2026).
			\bibitem{AppendixL} A. V. Yakushev, Anhang L: Fraktale Koordinationsfehlerskalierung, in \emph{YUCT} (2026).
			\bibitem{AppendixC} A. V. Yakushev, Anhang C: Bindungsenergien und das 0.67-Gesetz, in \emph{YUCT} (2026).
			\bibitem{AppendixG} A. V. Yakushev, Anhang G: Quantengravitation und Koordination, in \emph{YUCT} (2026).
			\bibitem{Schenk2004} P. Schenk et al., \emph{Icarus} \textbf{168}, 352–370 (2004).
			\bibitem{Dohnanyi1969} J. S. Dohnanyi, \emph{J. Geophys. Res.} \textbf{74}, 2531–2554 (1969).
			\bibitem{Tanaka2001} H. Tanaka, K. Ohtsuki, \emph{Icarus} \textbf{149}, 210–222 (2001).
			\bibitem{USGSNEIC} USGS National Earthquake Information Center, Erdbebenkatalog, \url{https://earthquake.usgs.gov/earthquakes/search/}.
			\bibitem{GutenbergRichter1954} B. Gutenberg, C. F. Richter, \emph{Seismicity of the Earth and Associated Phenomena}, Princeton University Press (1954).
			\bibitem{Aki1965} K. Aki, \emph{Bull. Earthq. Res. Inst.} \textbf{43}, 237–239 (1965).
			\bibitem{Stauber2015} J. M. Stauber et al., \emph{Langmuir} \textbf{31}, 2353–2360 (2015).
			\bibitem{Picknally1977} L. C. Picknally, \emph{J. Colloid Interface Sci.} \textbf{59}, 240–250 (1977).
			\bibitem{WebPlotDigitizer} A. Rohatgi, WebPlotDigitizer, \url{https://automeris.io/WebPlotDigitizer}.
		\end{thebibliography}
		
	\end{multicols}
\end{document}